\chapter{Occupation states and two Mellin actions}

An energy proportional to $\log m$ can be represented by a point
mass at $m^{-1}$. The Mellin transform of that mass is a Dirichlet
monomial. This gives an exact map from occupation states to
distributions. Its behavior at a zero of an entire function can
therefore be calculated before choosing a Hilbert norm on the
residue values.

All vector spaces are over $\mathbb C$. Inner products are linear
in the first variable. The coordinate $x>0$ is dimensionless,
and $t=\log x$ uses the real logarithm. Put
$\mathscr D=C_c^\infty(\mathbb R_{>0})$. Its elements are smooth
functions supported in a compact subset of the open half-line.
For $f\in\mathscr D$, define
\[
 F(s)=\mathcal Mf(s)=\int_0^\infty f(x)x^{s-1}\,dx.
\]
Compact support away from zero makes $F$ entire and permits every
spectral derivative under the integral. Define two operators on
this common invariant domain by
\[
 Qf=-(\log x)f,\qquad Df=x f'(x).
\]
Differentiation under the integral and integration by parts give
\begin{equation}
 \mathcal M(Qf)=-F',\qquad \mathcal M(Df)=-sF.
 \label{mo:two-actions}
\end{equation}
The boundary term $[f(x)x^s]_0^\infty$ is zero. Thus $Q$ and $D$
already have different actions on entire Mellin transforms.

\section{An exact occupation embedding}

Let $\mathscr E=c_{00}(\mathbb N_{\geq1})$ be the space of finite
sequences, with basis $v_m$ and inner product
$\langle a,b\rangle=\sum_m a_m\overline{b_m}$. Define
\[
 Nv_m=(\log m)v_m.
\]
Its Hilbert completion is $\ell^2(\mathbb N_{\geq1})$. The maximal
self-adjoint extension of $N$ has domain
\[
 \operatorname{Dom}N=
 \left\{a:\sum_m(\log m)^2|a_m|^2<\infty\right\}.
\]
Indeed finite truncations converge in the graph norm, and the
adjoint equation on the coordinate vectors imposes exactly this
domain. The coefficients $(\log m\pm i)^{-1}$ give everywhere
bounded inverses for $N\pm i$.

A distribution acts complex-linearly on smooth test functions.
The Dirac distribution $\delta_a$ satisfies
$\langle\delta_a,\phi\rangle=\phi(a)$; this convention uses the
additive measure $dx$. A compactly supported distribution $u$
acts on $x^{s-1}$ after insertion of a smooth cutoff equal to one
near its support. The result is independent of the cutoff. This
defines its entire Mellin transform.

\begin{proposition}[occupation distributions]
\label{mo:occupation}
The map
\begin{equation}
 T:\mathscr E\longrightarrow\mathscr E'(\mathbb R_{>0}),
 \qquad Tv_m=m^{-1}\delta_{m^{-1}}
 \label{mo:embedding}
\end{equation}
is injective. Here $\mathscr E'$ denotes compactly supported
distributions. Multiplication of distributions by the smooth
function $-\log x$ satisfies
\[
 QT=TN,\qquad
 \mathcal M(Ta)(s)=\sum_m a_m m^{-s}.
\]
\end{proposition}
\begin{proof}
A bump supported near $m^{-1}$, excluding the finitely many other
support points, recovers $a_m/m$ from $Ta$. This proves injectivity.
On that point mass, $-\log x$ takes the value $\log m$.
Finally,
$m^{-1}(m^{-1})^{s-1}=m^{-s}$, proving the Mellin formula.
\end{proof}

The factor $m^{-1}$ is fixed by the additive-distribution convention.
If point masses were instead introduced relative to $dx/x$, their
conversion to additive distributions would contain this same factor.
The weighted occupation lattice with coordinate vectors $d_m$ of
squared norm $m^{-1}$ maps isometrically to the occupation space by
$d_m\mapsto m^{-1/2}v_m$. Its resulting distribution is
$m^{-3/2}\delta_{m^{-1}}$, with Mellin transform $m^{-s-1/2}$.

\begin{proposition}[extension and the analytic half-plane]
\label{mo:dirichlet-domain}
The same locally finite sum defines an injective map
$T:\ell^2\to\mathscr D'(\mathbb R_{>0})$. It is continuous for
the strong distribution topology, and $QT=TN$ on
$\operatorname{Dom}N$. For $a\in\ell^2$, the Dirichlet series
\[
 F_a(s)=\sum_m a_m m^{-s}
\]
and all its spectral derivatives converge absolutely and locally
uniformly on $\operatorname{Re}s>1/2$. No continuation across the
line $\operatorname{Re}s=1/2$ is part of this construction.
\end{proposition}
\begin{proof}
A compact subset of $\mathbb R_{>0}$ meets only finitely many points
$m^{-1}$. Thus the sum acts on every test function. Bounded sets
of tests have a common compact support and uniform bounds on all
derivatives. Only a fixed finite set of coefficients occurs on such
a set, so Cauchy--Schwarz proves strong continuity. Separate bumps
still recover every coefficient, and multiplication proves the
operator identity.

For $\sigma>1/2$ and integer $k\geq0$,
\[
 \sum_m |a_m|(\log m)^k m^{-\sigma}
 \leq\|a\|_2
 \left(\sum_m(\log m)^{2k}m^{-2\sigma}\right)^{1/2}<\infty.
\]
Integral comparison proves convergence, uniformly when $\sigma$
stays above $1/2+\epsilon$. Termwise differentiation is therefore
valid. The estimate supplies no holomorphic neighborhood of a
point on the boundary line.
\end{proof}

For a distribution $u$, set $\langle u',\phi\rangle=-\langle
u,\phi'\rangle$ and $Du=xu'$. The same calculation gives
$\mathcal M(Du)=-s\mathcal Mu$ for compactly supported $u$.
In particular,
\[
 D(Tv_m)=m^{-1}\bigl(m^{-1}\delta'_{m^{-1}}-\delta_{m^{-1}}\bigr).
\]
Dilation does not preserve the span of the point masses in
\eqref{mo:embedding}. The identity $QT=TN$ cannot be replaced by
an identity involving $D$.

\section{Residues of occupation distributions}

Let $\Xi$ be a nonzero entire function. Fix a nonempty finite set
$Z$ of distinct zeros. For $\rho\in Z$, let $r_\rho\geq1$ be its
multiplicity, and write
\[
 \Xi(s)^{-1}=
 \sum_{\ell=-r_\rho}^{\infty}b_{\rho,\ell}(s-\rho)^\ell,
 \qquad b_{\rho,-r_\rho}\ne0.
\]
For an entire $F$, define its residue coordinates by
\begin{equation}
 R_{\rho,k}F=\operatorname*{Res}_{s=\rho}
 \frac{(s-\rho)^kF(s)}{\Xi(s)},
 \quad 0\leq k<r_\rho.
 \label{mo:jets}
\end{equation}
The residue is the coefficient of $(s-\rho)^{-1}$. These coordinates
record the full principal part of $F/\Xi$. Direct multiplication
of the Taylor and Laurent series gives
\begin{equation}
 R_{\rho,k}F=
 \sum_{j=0}^{r_\rho-k-1}
 b_{\rho,-k-j-1}\frac{F^{(j)}(\rho)}{j!}.
 \label{mo:jet-formula}
\end{equation}
The transformation from the Taylor jet of order $r_\rho-1$ to
these coordinates is invertible, since its successive leading
coefficients are $b_{\rho,-r_\rho}$.

\begin{proposition}[the occupation residue defect]
\label{mo:occupation-defect}
Define $A_0a=(R_{\rho,0}F_a)_{\rho\in Z}$ on $\mathscr E$.
Let $K_\lambda$ multiply the $\rho$ coordinate by a prescribed
number $\lambda_\rho$. Then $A_0N=K_\lambda A_0$ fails on
$\mathscr E$ for every choice of the numbers $\lambda_\rho$.
At a simple zero,
\begin{equation}
 (A_0v_m)_\rho=\frac{m^{-\rho}}{\Xi'(\rho)},\qquad
 ((A_0N-K_\lambda A_0)v_m)_\rho
 =\frac{(\log m-\lambda_\rho)m^{-\rho}}{\Xi'(\rho)}.
 \label{mo:simple-occupation}
\end{equation}
\end{proposition}
\begin{proof}
Formula \eqref{mo:jet-formula} yields
\[
 (A_0v_m)_\rho=m^{-\rho}p_\rho(\log m),\qquad
 p_\rho(u)=\sum_{j=0}^{r_\rho-1}
 b_{\rho,-j-1}\frac{(-u)^j}{j!}.
\]
The polynomial has degree $r_\rho-1$ and nonzero leading
coefficient. Among $m=1,\ldots,r_\rho+1$, at least two values
$p_\rho(\log m)$ are nonzero. The proposed identity would make
$\lambda_\rho$ equal to two distinct logarithms. At a simple zero,
$m=1$ and $m=2$ already decide the question.
\end{proof}

For the local double-zero model $\Xi(s)=(s-\rho)^2$,
\[
 R_{\rho,0}(m^{-s})=-(\log m)m^{-\rho},\qquad
 R_{\rho,1}(m^{-s})=m^{-\rho}.
\]
The scalar residue kills $v_1$, but $v_2$ and $v_3$ still contradict
a fixed diagonal energy. This local entire function is a model
of the multiplicity calculation, not an identification with an
arithmetic entire function.

For general $a\in\mathscr E$, equation \eqref{mo:two-actions} gives
\[
 (A_0Na)_\rho=\operatorname*{Res}_{s=\rho}
                      \frac{-F_a'(s)}{\Xi(s)}.
\]
Thus allowing point distributions repairs the smooth-support
obstruction to $T$ itself. It leaves the derivative in the residue
composition. The target $\lambda_\rho=\operatorname{Im}\rho$ has
exactly the defect in \eqref{mo:simple-occupation}.

\section{A one-frequency obstruction before completion}

For real $c$, put
\[
 \mathscr H_c=L^2(\mathbb R_{>0},x^{2c-1}\,dx),\qquad
 (U_cf)(t)=e^{ct}f(e^t).
\]
Change of variables gives $\|U_cf\|_{L^2(dt)}=\|f\|_{\mathscr H_c}$.
Fix a real frequency $\gamma_0$ and choose
$\psi\in C_c^\infty(\mathbb R)$ with $\int\psi(t)\,dt=1$.
For $L>0$, define
\[
 g_L(t)=L^{-1}\psi(t/L)e^{-i\gamma_0t},\qquad f_L=U_c^{-1}g_L.
\]
Then $f_L\in\mathscr D$ and
\begin{equation}
 \|f_L\|_{\mathscr H_c}^2=L^{-1}\|\psi\|_2^2\longrightarrow0,
 \qquad \mathcal Mf_L(c+i\gamma_0)=1.
 \label{mo:continuous-sequence}
\end{equation}
For a simple zero $\rho=c+i\gamma_0$, its residue coordinate
therefore remains $1/\Xi'(\rho)$ while the source norm tends to zero.
A linear operator is \emph{closable} if every sequence of domain
vectors tending to zero whose images converge has image limit zero.
This necessary and sufficient condition says that the closure of
its graph is still the graph of a function. Equation
\eqref{mo:continuous-sequence} proves that this one-coordinate
evaluation is not closable on $\mathscr H_c$.

There is an equally explicit occupation sequence. Let
$H_M=\sum_{m=1}^M m^{-1}$ and fix $\rho=1/2+i\gamma_0$. Set
\[
 a_m^{(M)}=
 \begin{cases}H_M^{-1}m^{-1/2+i\gamma_0},&m\leq M,\\0,&m>M.
 \end{cases}
\]
Then
\begin{equation}
 \|a^{(M)}\|_{\ell^2}^2=H_M^{-1}\longrightarrow0,
 \qquad F_{a^{(M)}}(\rho)=1.
 \label{mo:atomic-sequence}
\end{equation}
Divergence of $H_M$ follows from comparison with $\int_1^{M+1}dx/x$.
Consequently one critical-line evaluation, and hence its nonzero
simple-residue multiple as a scalar operator, is not closable on
occupation $\ell^2$. This statement concerns that scalar operator.
It does not infer nonclosability of an arbitrary vector-valued
map merely from one unbounded coordinate.

\begin{proposition}[simultaneous critical-line occupation evaluations]
\label{mo:finite-occupation-graph}
Let $\rho_j=1/2+i\gamma_j$, $1\leq j\leq n$, have distinct
real ordinates. The map
$a\mapsto(F_a(\rho_j))_{j=1}^n$ on $\mathscr E$ has dense graph
in $\ell^2\oplus\mathbb C^n$. The same holds after multiplication
of its coordinates by any fixed nonzero constants, including
simple-residue factors.
\end{proposition}
\begin{proof}
Let $S_M$ be the matrix with entries $(S_M)_{jm}=m^{-\rho_j}$,
$1\leq m\leq M$. Its Gram matrix $G_M=S_MS_M^*$ has diagonal
entries $H_M$. For $j\ne k$, its entry is
$\sum_{m\leq M}m^{-1+i(\gamma_k-\gamma_j)}$, which is bounded
uniformly in $M$. Indeed its integral is bounded by
$2/|\gamma_k-\gamma_j|$, and the difference between the sum and
integral is bounded by the integral of the absolute derivative,
which is a constant times $\int_1^\infty x^{-2}\,dx$.
Thus $G_M=H_MI+O(1)$ in the fixed matrix dimension. For large $M$
it is invertible and $\|G_M^{-1}\|\to0$.
Given $y\in\mathbb C^n$, extend
$a^{(M)}=S_M^*G_M^{-1}y$ by zero outside $m\leq M$. Then
\[
 S_Ma^{(M)}=y,\qquad
 \|a^{(M)}\|_2^2=y^*G_M^{-1}y\longrightarrow0.
\]
The graph closure contains every $(0,y)$. Subtracting such
vectors from the graph over the dense finite-sequence core proves
density in the entire product.
\end{proof}

\begin{proposition}[finite occupation jets]
\label{mo:occupation-jet-graph}
For finitely many distinct zeros on $\operatorname{Re}s=1/2$,
the full residue map on $\mathscr E$ has dense graph in
$\ell^2\oplus\mathscr J_Z$. Its range is all of $\mathscr J_Z$.
For any nonempty finite zero set, its kernel on $\mathscr E$
is not invariant under $N$.
\end{proposition}
\begin{proof}
A linear combination of the Taylor evaluation kernels is
$m^{-1/2}q(\log m)$, where
$q(t)=\sum_jp_j(t)e^{-i\gamma_jt}$. For nonzero coefficients,
independence in Lemma~\ref{mo:exponential-polynomials} gives $q\ne0$.
For $h(x)=x^{-1}|q(\log x)|^2$, the derivative is bounded by
$Cx^{-2}(1+\log x)^d$ for some integer $d$.
This is integrable on $[1,\infty)$. Hence the difference between
$\sum_{m\leq M}h(m)$ and $\int_1^M h(x)\,dx$ stays bounded.
The integral equals $\int_0^{\log M}|q(t)|^2dt$ and tends to
infinity by the positive-half-line estimate in that lemma.
No nonzero combination of the kernels belongs to $\ell^2$.
The orthogonal-complement argument from Theorem~\ref{mo:finite-graph}
now proves dense graph. A linear subspace of a finite-dimensional
space is closed, so density of the range proves surjectivity.

For the last assertion, suppose the kernel were $N$-invariant.
There would be an induced linear operator on the finite-dimensional
range of the residue map. The image of $v_m$ is nonzero: its
last coordinate at any zero is
$b_{\rho,-r_\rho}m^{-\rho}\ne0$.
That image would be an eigenvector with eigenvalue $\log m$.
Infinitely many distinct eigenvalues are impossible in finite
dimension. This contradiction applies without zero-placement
or simplicity assumptions.
\end{proof}

\section{The Mellin-diagonal dilation generator}

The norm of $\mathscr H_c$ dictates the dilation amplitude. Define
\[
 (V_uf)(x)=e^{cu}f(e^u x),\qquad u\in\mathbb R.
\]
Substitution $y=e^u x$ proves $\|V_uf\|_{\mathscr H_c}=
\|f\|_{\mathscr H_c}$. In logarithmic coordinates, $V_u$ is
translation $g(t)\mapsto g(t+u)$. Its infinitesimal operator on
$\mathscr D$ is $D+c$. With the convention $V_u=e^{-iuP_c}$,
\begin{equation}
 P_c=i(D+c),\qquad
 \mathcal M(P_cf)(s)=i(c-s)F(s).
 \label{mo:dilation}
\end{equation}
In particular $\mathscr H_{1/2}=L^2(dx)$ and
$P_{1/2}=i(x\partial_x+1/2)$. The half-shift comes from the
measure, and the factor $i$ comes from the stated unitary-group
convention. On $L^2(dx/x)$ one has $c=0$ and no half-shift.

\begin{lemma}[the normalized Mellin unitary]
\label{mo:mellin-unitary}
The map
\[
 (\mathcal U_cf)(\tau)=\frac1{\sqrt{2\pi}}
                          \mathcal Mf(c+i\tau)
\]
on $\mathscr D$ extends to a unitary
$\mathscr H_c\to L^2(\mathbb R,d\tau)$. The closure of $P_c$ is
self-adjoint and is multiplication by $\tau$ under this unitary.
Its domain is
\[
 \left\{f\in\mathscr H_c:
       \tau\mathcal U_cf(\tau)\in L^2(d\tau)\right\}.
\]
Equivalently, $U_cf$ has weak derivative in $L^2(dt)$.
\end{lemma}
\begin{proof}
For $g=U_cf$, the displayed map is the Fourier transform
$\widehat g(\tau)=(2\pi)^{-1/2}\int g(t)e^{i\tau t}\,dt$.
Its isometry can be proved directly. For Schwartz functions
$g,h$ and $\epsilon>0$, Fubini and the Gaussian integral give
\[
 \int\widehat g(\tau)\overline{\widehat h(\tau)}
                      e^{-\epsilon\tau^2}\,d\tau
 =\iint g(t)\overline{h(u)}
 \frac{e^{-(t-u)^2/(4\epsilon)}}{2\sqrt{\pi\epsilon}}\,dt\,du.
\]
The Gaussian kernels have integral one and converge as an
approximate identity. Splitting their integrals into
$|t-u|<\delta$ and its complement, then using translation
continuity in $L^2$, proves convergence to $\langle g,h\rangle$.
For $h=g$, monotone convergence proves equality of the two squared
norms. Polarization gives the inner-product identity.
Integration by parts shows that Fourier transforms of Schwartz
functions are Schwartz. The same Gaussian calculation gives
Fourier inversion by first inserting $e^{-\epsilon\tau^2}$ and
then letting $\epsilon$ tend to zero. Thus the range contains the
Schwartz space. Its density and the isometry give a unitary on
$L^2$. Cutoffs and convolution by smooth approximate identities
prove the required density of compact smooth functions.

On compact smooth functions, $\widehat{ig'}=\tau\widehat g$.
Multiplication by $\tau$ on its maximal domain is self-adjoint:
testing against compactly supported functions identifies its
adjoint domain, and division by $\tau\pm i$ gives its resolvents.
The inverse Fourier transform identifies this domain with weakly
differentiable $L^2$ functions whose derivative is in $L^2$.
Such functions can be approximated in the norm
$\|g\|_2+\|g'\|_2$ by cutoffs followed by smooth convolution.
Indeed the cutoff derivative term is bounded by
$C\|g\|_2/L$ for a cutoff varying at scale $L$.
Therefore compact smooth functions form a graph core. This proves
the closure and domain assertions.
\end{proof}

The spectrum of $P_c$ is all of $\mathbb R$ and has no eigenvectors.
A nonzero $L^2$ function cannot be supported at one frequency.
Functions supported in shrinking intervals around any prescribed
frequency give approximate eigenvectors. In contrast, the
occupation operator has the eigenbasis $v_m$ and eigenvalues
$\log m$.

In logarithmic coordinates, $Q=-t$ and $P_c=i\partial_t$.
On compact smooth functions,
\[
 [P_c,Q]=-i\,\mathrm{id}.
\]
With $\hbar=1$, these are a position observable with reversed
orientation and its conjugate momentum. The distribution $Tv_m$
is localized at $t=-\log m$. Its Mellin transform is a plane
wave multiplied by $m^{-c}$ on the line $\operatorname{Re}s=c$;
it is not a frequency point mass. This distinction accounts for
both actions in \eqref{mo:two-actions}.

\section{Dilation on full residue jets}

Let $\mathscr J_Z=\bigoplus_{\rho\in Z}\mathbb C^{r_\rho}$,
with coordinates $a_{\rho,k}$ as in \eqref{mo:jets}. Set
$a_{\rho,r_\rho}=0$ and define
\begin{equation}
 (K_ca)_{\rho,k}=z_{\rho,c}a_{\rho,k}-i a_{\rho,k+1},
 \qquad z_{\rho,c}=i(c-\rho).
 \label{mo:jet-generator}
\end{equation}
Write $\mathcal R_Zf=(R_{\rho,k}\mathcal Mf)_{\rho,k}$.

\begin{proposition}[the full jet comparison]
\label{mo:jet-comparison}
On $\mathscr D$,
\[
 \mathcal R_ZP_c=K_c\mathcal R_Z.
\]
At a simple zero $\rho=\sigma+i\gamma$, the multiplier is
$z_{\rho,c}=\gamma+i(c-\sigma)$. It equals the ordinate
$\gamma$ exactly when $\sigma=c$. At a zero of multiplicity
$r$, the full block is a Jordan block with eigenvalue $z_{\rho,c}$
and superdiagonal $-i$.
\end{proposition}
\begin{proof}
In \eqref{mo:dilation}, write
$i(c-s)=i(c-\rho)-i(s-\rho)$ and take the residue after
multiplication by $(s-\rho)^k$. The second term is the next
coordinate in \eqref{mo:jets}; it vanishes at the last coordinate.
The superdiagonal is nonzero, so the nilpotent part has exact
order $r$.
\end{proof}

For a double zero, the scalar residue alone obeys
\[
 R_{\rho,0}\mathcal M(P_cf)=
 z_{\rho,c}R_{\rho,0}F-iR_{\rho,1}F.
\]
Dropping the second coordinate drops an actual term. Repeating
$F(\rho)$ in two coordinates instead gives a diagonal action, but
its image has dimension one. It does not retain the full
multiplicity-two jet. Nor does the jet comparison repair the
occupation map: its last coordinate on $Tv_m$ is
$b_{\rho,-r_\rho}m^{-\rho}\ne0$. Intertwining $N$ with $K_c$
would force $\log m=z_{\rho,c}$ for every $m$.
